\subsection{Closure-flow curvature regime}

\begin{proposition}[Directed joins are the Scott-continuity test carrier]\label{thm:directed-family-carrier}\lean{Obsmetrologia.Closure.Scott.DirectedFamily}\leanok\uses{thm:closure-regime-data, thm:kle}
A directed family is the finite charge-lattice projection of a directed subset.
It is the carrier on which a closure operator is tested for commuting with
directed joins.
\end{proposition}

\begin{theorem}[Finite charge closures are Scott-continuous]\label{thm:finite-scott-continuous}\lean{Obsmetrologia.Closure.Scott.finite_scott_continuous}\leanok\uses{thm:directed-family-carrier, thm:closure-operator, thm:spine-is-strict-regime}
On the finite charge carrier, every directed join is attained.  Therefore every
charge closure commutes with directed joins and sits on the strict/Scott side
of the regime projection.
\end{theorem}

\begin{proposition}[A non-Scott firing is failure to commute with a directed join]\label{thm:nsc-firing-carrier}\lean{Obsmetrologia.Flow.NSCFiring.IsNSC}\leanok\uses{thm:finite-scott-continuous}
An NSC firing is the precise obstruction equation: firing after the directed
join is not the directed join of the firings.  This is the incidence point
where infinitary antecedents enter the proof tree.
\end{proposition}

\begin{theorem}[Scott-continuity is absence of NSC firing]\label{thm:scott-continuous-iff-no-nsc-firing}\lean{Obsmetrologia.Flow.NSCFiring.scott_continuous_iff_no_nsc_firing}\leanok\uses{thm:nsc-firing-carrier, thm:finite-scott-continuous}
A charge closure is Scott-continuous exactly when no directed-family prefix is
an NSC firing for its operator.  This turns the regime flag into an incidence
condition.
\end{theorem}

\begin{theorem}[Finite-carrier NSC firing facts]\label{thm:nsc-firing-source-of-aperiodic}\lean{Obsmetrologia.Flow.NSCFiring.no_closure_nsc_firing_and_forallAnte_witness}\leanok\uses{thm:scott-continuous-iff-no-nsc-firing, thm:regime-cohomology-vanishes-iff}
Genuine finite charge closures have no NSC firing, while a concrete
forall-antecedent firing operator witnesses one.  The statement says nothing about
regimes: that NSC firings source the non-Scott regimes is open, since the regime
datum of the module is chosen by hand and not computed from a firing.
\end{theorem}

\begin{theorem}[Transfinite closure flow has no settling-limit blowup]\label{thm:transfinite-closure-flow-no-blowup}\lean{Obsmetrologia.Valuation.ModalFrame.transfinite_closure_flow_no_blowup}\leanok\uses{thm:scott-continuous-iff-no-nsc-firing, thm:finite-scott-continuous}
For the modal-frame closure flow indexed by the internal ordinal, every
settling limit sequence is continuous and has no NSC firing.  The theorem
therefore gives the discrete global-regularity statement that is currently
proved: no limit-stage blowup occurs once the predecessor chain has reached the
settled value.
\end{theorem}

\begin{theorem}[Raw all-limit regularity is a frontier obligation]\label{thm:raw-all-limit-regularity-frontier}\notready\uses{thm:transfinite-closure-flow-no-blowup}
The stronger slogan "continuity at every raw limit constructor" is not yet a
proved theorem.  The current Brouwer ordinal admits degenerate fundamental sequences,
so this claim needs an additional carrier for genuine/cofinal limit
presentations before it can be promoted.
\end{theorem}

\begin{theorem}[Curvature labels on the declared regime table]\label{thm:curvature-regime-correspondence}\lean{Obsmetrologia.Flow.CurvatureRegime.curvature_regime_table}\leanok\uses{thm:nsc-firing-source-of-aperiodic, thm:regime-cohomology-vanishes-iff}
The labels ``declared flat'' and ``declared curved'' split the four regimes, and the
declared cohomology profile (a table; no \(H^1\) is computed) agrees with them.  This holds
by the definitions.  That the regime classification is a classification by curvature, with
\(H^1\) the curvature module, is open.
\end{theorem}

\begin{theorem}[The curved row beside an NSC firing witness]\label{thm:curvature-rank-has-nsc-firing}\lean{Obsmetrologia.Flow.CurvatureRegime.declaredCurved_row_beside_nsc_witness}\leanok\uses{thm:curvature-regime-correspondence}
The curved row of the declared tables is stated beside an actual NSC firing witness of the
forall-antecedent operator.  No map links the witness to the row, so
\(\operatorname{rank} H^1 = |L_{ns}|\) is open.
\end{theorem}
